% 原书第 iii 页（PDF 第 7 页）
\BookPage{iii}{7}
\BookHeading{序}{序}{}
我对陈捷先教授心仪很久，却相见恨晚。

上世纪七十年代，海峡阻隔，资讯不通。先生已经出版《满洲丛考》和《清史杂笔》等多篇论文和多部专着，我却对先生的学术研究资讯一无所知。

上世纪八十年代，我作为清史访问学者到了美国。在进行学术访问和交流期间，我在图书馆的书架上，在同友朋的交谈中，才看到或听到「陈捷先」这个名字，也才看到《满洲丛考》这部书。那个时候，我的重点研究领域是清朝开国史，所以对陈先生的学术成果和学术活动格外关注。

1989年我倡议建立北京社会科学院满学研究所，接着筹备「首届国际满学研讨会」。这期间，我同陈先生不断有通信、传真和电话联系。1992年，我在北京前门饭店主持召开「首届国际满学研讨会」，陈先生当时正在筹备第三十五届世界阿尔泰学会会议，难以分身，但还是寄来〈三田渡满文清太宗功德碑研究〉的论文。随后我在《满学研究》第一辑，刊出陈先生的这篇论文，并在「满学家」专栏里，载文介绍陈捷先教授和神田信夫教授两位国际满学界的大家。然而，我和陈先生仍没机会见面。

直到同年秋后，应陈捷先教授邀请，我同王钟翰、韦庆远、王戎笙、冯尔康、徐艺圃、刘耿生、林岷、柏桦一行，作为大陆第一批社会科学工作者，登上美丽的台湾岛，出席陈捷先教授主持的「海峡两岸清史档案学术研讨会」，我和陈先生终于有了第一次握手，第一次拥抱。

从此开了先河，我后来多次赴台交流。每念及此，感激捷公。

性情率真，袒露肺腑，是陈先生给我留下的第一个印象。我因长期受着无形绳索的捆绑、政治运动的拷打，夹着尾巴做人，谨小慎微，万分拘谨。到台湾第一个感受是，陈先生等学人开朗豪放，言所欲言，纵情畅饮，通宵达旦。


% 原书第 iv 页（PDF 第 8 页）
\BookPage{iv}{8}
治学认真，著作等身，是陈先生给我留下的又一个印象。陈先生早在1956年就毕业于台湾大学历史学系，1959年获台湾大学历史研究所硕士学位，后应邀到美国加入哈佛大学访问学者研究。先生六十年如一日，勤奋研究，笔耕不辍。《满文清实录研究》、《满文清本纪研究》、ManchuArchival Materials（满文档案资料概述），近年出版《蒋良骐《东华录》研究》、《清史论集》，以及清十二朝的《清史事典》，还有《努尔哈齐写真》、《皇太极写真》、《顺治写真》、《康熙写真》、《雍正写真》和《乾隆写真》等，尤其是后者，一人之力，连出六册。

满腔热忱，肝胆相照，是陈先生给我留下的再一个印象。陈先生于我，可谓亦师亦友，多次来京，欢快相聚，倾心交谈。我也多次受邀飞抵台湾，出席学术活动。屈指算来，已有八次。记得2004年10月20日，陈捷先教授提出两岸学者共同编着出版《清代与台湾》丛书的重要倡议，同我和冯尔康先生商量，三人一拍即合。对于《清代台湾》一书，陈先生不仅主持全局，拟定编写体例，通审全稿，并撰写为全书之纲的〈通纪〉篇，概述清朝以前的台湾历史，重点阐述清代台湾的历史与文化，特别叙述清代台湾与中央政府的关系。这部书稿，历时五年，多方协作，终于由九州出版社出版。

2009年，应台湾故宫博物院周功鑫院长、冯明珠副院长邀请，我到台北出席「两岸故宫第一届学术研讨会—为君难：雍正其人其事及其时代」。会间，我和冯尔康先生又发现了新的缘分一原来我们俩是同年同月同日生。在高兴之际，同为扬州籍的两位清史才子一陈先生和冯先生共商于2010年在扬州举办「盛清社会与扬州」学术研讨会。这次学术盛会，得到扬州市王燕文书记、谢正义市长等领导支持，按期举行，圆满结束。会间，又套出一个主意，就是以此次会议论文为基础，再约请未及与会的先生撰写文章，合作出版一本论文集。事情也巧，一算年份，恰是陈捷先教授八十华诞，大家相约以此论集作为向陈捷先教授八十大寿敬示的一个学术祝贺。

现在，《盛清社会与扬州研究》文集即将出版，谨遵操持此事的冯明珠副院长之嘱，撰写上文，忝做小序。

\BookSignature{阎崇年于北京四合书屋2011年8月28日}


% 原书第 v 页（PDF 第 9 页）
\BookPage{v}{9}
\BookHeading{捷公祝寿集序言}{捷公祝寿集序言}{}
人云韶华只为少年留，吾曰韶华岂独在少年！八秩寿翁陈捷公，活跃两岸枫叶国，辉煌学术再创造，展现人生新境域。韶华为捷公永驻，观念为捷公更新；友朋给捷公拜寿，预祝未来更绚丽。

1995年捷公在台湾大学荣退，没有隐居于枫叶之国加拿大，而是穿梭于海峡两岸，继续学术研究事业，进行三方面的活动，一为潜心写作，二是推动学术交流，三是办学。荣退后捷公问世新着之多，令人惊讶不已！有面向大众读者的《努尔哈齐写真》、《皇太极写真》、《顺治写真》、《康熙写真》、《雍正写真》、《乾隆写真》、《慈禧写真》及《宣统写真》八部，由远流出版公司、商务印书馆、紫禁城出版社、浙江文艺出版社多家书社印行，风行于两岸广大读者中。学术专着《清代台湾方志研究》、《东亚古方志学》和《蒋良骐及其《东华录》研究》多部连袂问世。关注现实社会，撰着「以古鉴今，痛下针砭」的学术杂文集《不剃头与两国论》。主编《清史事典》一套，计十二册；与阎崇年教授共同主编《清代台湾研究丛书》，合著《清代台湾》；合作主编《清史论集》。创办宜兰佛光学院历史学系所，授业解惑。2003年在宜兰主办第一届国际清史研讨会；2010年与阎崇年教授、尔康共同发起召开「盛清扬州社会高端论坛」研讨会。古稀之年，贡献给社会几十部著述，青壮年学者亦难于做到。人们说老年人做点事是发挥余热，是「老有所为」，在捷公这里就不准确了，不是余辉余热、有所为，而是老当益壮、大有作为，发出学术光芒，惠及读者大众和社会。捷公的学术青春常在，韶华永驻，颇为符合人文学科的研究状态。明末清初三大家顾炎武、黄宗羲、王夫之都在古稀耄耋之年，笔耕不辍，如果他们中年谢世，或者停止钻研，他们就很难成为「三大家」了。人到老年，有知识积累、人生阅历的治学优势，有除却一些杂务干扰的有利条件，便于学术总结，令研究更上一层楼，这大约是人文学科研究的规律，捷公学术上的永葆韶华，是为生动有力的证明。


% 原书第 vi 页（PDF 第 10 页）
\BookPage{vi}{10}
往前说的话，我认识的捷公，是清史大家，是文献学家，是推动世界汉学研究的文化使者，是颇有「叱咤风云」味道的权威学者。我不想在这里罗列捷公著述的书单，评论他在清史、方志学、谱牒学研究方面的学术贡献，倒想绍述他推动学术研究的业绩。捷公的好友张存武教授在〈云霞秋色忆华年〉文中，说捷公「活动力强，交际人缘广」，「深谋远虑，做参谋总长最适合」。可知捷公虽不做官，但组织能力很强。这正是他的优势，他用来组织学术活动，举办国际间的学术会议。由捷公创设和主持召开的会议有：「亚洲族谱学术研讨会」、「中韩历史关系国际学术会议」、「中国域外汉籍国际学术会议」、「中琉历史关系国际学术会议」等等，这些学术会议都是一而再再而三地多次进行，比如族谱学研讨会举行九次，第八届是捷公与香港大学赵令扬教授联合举办的，第九届是和我所服务的南开大学等单位合作的。中琉关系史研讨会持续进行，台北、琉球、北京、福州等地轮流坐庄。捷公开创的这些研究活动，促进谱牒学、方志学、韩国学、琉球学、汉学诸多领域的研究蓬勃发展，培养了一批中外青年学者，是以人们说捷公「既开风气且为师」。写到这里，我想起捷公于1992年召开的第三十五届世界阿尔泰学会会议，阿尔泰学基本上是关于东方的学问，由于历史的原因，会议从来是在西方举行的，这一届的会议欧洲某国要求主办，捷公与其竞争主办权，终因学术各方面的实力，如愿以偿，会议在台北圆山饭店揭幕。此次会议，承蒙捷公盛情，邀约九位大陆学者出席，我遂得首次观光台湾。捷公举办这一系列的学术会议，有着强烈的传播中华文化的使命感，他是要让东方人、西方人了解历史上中国与东北亚、东南亚的密切关系及中华文化的远播，让世界学人了解中国和日本、韩国、越南学者的研究状况和学术观点。事实表明捷公功在中华文化的向世界传播，功在中华文化的继承发扬，功在两岸的学术交流。他是文化交流使者、推动者，正是在这个意义上我说他有叱咤风云的学术气概，令人佩服。

说到两岸史学界的交往，阿尔泰研讨会大陆学者九人赴台，人数之多是前所未有的，此乃捷公开风气之先的举措。两岸史学工作者切磋交流，只有在改革开放之后才有可能，第一次提供大规模交流平台的是前述港大教授赵令扬，他于1985年举办国际明清史研讨会，两岸各有近二十名学者出席，我就是在这次会上认识了捷公和其他台湾学友。所以我以为两岸三地的史学家交往，是赵令扬教授在香港、捷公在台北先后创造了平台，功在史学界。


% 原书第 vii 页（PDF 第 11 页）
\BookPage{vii}{11}
捷公为人豪爽，饮酒中尽情地展现出来。我第一次感受深切的是在阿尔泰会议的晚宴上，一二百人的宴会，捷公挨桌敬酒，到大陆同行的坐席，只见他白酒、红酒、啤酒样样来得，年轻时颇能豪饮的王钟翰教授宅心仁厚，见此情状，劝他多吃菜，意在解酒勿醉，捷公表示领会他的好意，真是同胞之深情啊！捷公夙有酒仙之名，张存武教授说他喝酒不误事，是以酒会友。实在是挚友真知。捷公不仅是以文会友，更是以酒会友。这是中国文人的一种传统！

捷公是南开大学客座教授，对南开贡献良多。1996年我在台北，因捷公的关系，联经出版公司总经理姚为民先生宴请我们，无意中说到奖学金，豪侠好义的姚先生当即表示在南开设立奖学金，返津后我就和我的同事操办起来，每年给南开历史系学生评定发放姚为民奖学金，因姚先生是山东平度人，遵他嘱咐每年另给平度籍学生一人发奖，此事持续到我退休。1997年，捷公与彭炳进教授应邀访问南开，由捷公倡议，彭教授在南开历史系设立讲座基金，邀请京津学者讲学，出版演讲论文集，连续进行了三年。我和捷公是扬州同乡，史界同行，又都是少小离乡而乡情浓厚，这些因素将我们连结在一起，不过更主要的是我仰慕捷公的为人。上面说的那些话都是我的心里话呀！

2010年的扬州学术会议，两岸学者发表学术高见，提交宏文。经征得作者的同意，汇辑为庆祝捷公八秩华诞论文集，以文章祝贺学人，是最好的、最有价值的方式。这个集子采用中文繁体字，拙文则以简体字先期在大陆《安徽史学》今年第一期刊出，特在后记中表达对捷公祝贺：「…庆祝台湾大学荣退教授陈捷先先生八秩华诞，兹以简体文本披露，为挚友之庆。」参与文集编辑的大陆学者是阎崇年教授、常建华教授，捷公弟子龙贵居间联系，台湾学者冯明珠女士总其成。这是两岸学者为捷公祝寿，舍捷公难于有他人！

捷公，记得吧，前年（2009年），您的天津朋友诚挚邀请您们伉俪光临，您们何时成行呢？我们期待着！

谨祝捷公同夫人侯女士健康，愉悦！预祝捷公伉俪一如既往给大陆、台湾、加拿大航空业作「贡献」，令我们友朋时常相聚。

\BookSignature{冯尔康2011年8月22日于雪梨客室}
